\section*{Chapter \ref{ch:b3:solid-state} --- Solid-State Chemistry: Bands, Defects and Semiconductors}

\begin{solution}{exo:b3:solid-state:1}
$(E - \alpha)/|\beta| = -2\cos(k\pi/7)$: $-1.802$, $-1.247$, $-0.445$, $+0.445$, $+1.247$, $+1.802$. They span
$3.604|\beta|$, already 90\,\% of the band width $4|\beta|$ of the infinite chain.
\end{solution}

\begin{solution}{exo:b3:solid-state:2}
$2N$ ($N$ orbitals, two electrons each). Sodium, with one 3s electron per atom, half
fills its s band: a metal. Magnesium has two and would fill the s band exactly,
but the 3s and 3p bands overlap, so the highest occupied levels belong to a
partly filled combined band: a metal too.
\end{solution}

\begin{solution}{exo:b3:solid-state:3}
$\lambda \approx \qty{1239.8}{eV.nm}/E_g$: gallium phosphide \qty{549}{nm} (green); gallium nitride \qty{365}{nm}
(near ultraviolet). Blue diodes use gallium nitride alloyed with indium, whose
gap is smaller.
\end{solution}

\begin{solution}{exo:b3:solid-state:4}
$\ce{CaCl2} \xrightarrow{\ce{NaCl}} \mathrm{Ca_{Na}^{\bullet} + V_{Na}' + 2\,Cl_{Cl}^{\times}}$: two cation and two anion sites, as in NaCl; one Ca and
two Cl; charge $+1 - 1 = 0$. $\ce{Y2O3} \xrightarrow{\ce{ZrO2}} \mathrm{2\,Y_{Zr}' + 3\,O_O^{\times} + V_O^{\bullet\bullet}}$: two cation and four oxygen sites; two Y
and three O; charge $-2 + 2 = 0$.
\end{solution}

\begin{solution}{exo:b3:solid-state:5}
$T^{3/2} = 5196$ at \qty{300}{K}, $kT = \qty{0.02585}{eV}$. Silicon: $\sqrt{N_cN_v} = \qty{2.42e19}{cm^{-3}}$, $\eu^{-1.125/0.05170} = \num{3.55e-10}$,
$n_i = \qty{8.6e9}{cm^{-3}}$. Germanium: $\sqrt{N_cN_v} = \qty{7.16e18}{cm^{-3}}$, $\eu^{-0.661/0.05170} = \num{2.80e-6}$, $n_i = \qty{2.0e13}{cm^{-3}}$. Ratio
$\num{4.3e-4}$: the smaller gap of germanium gives it over two thousand times more
carriers.
\end{solution}

\begin{solution}{exo:b3:solid-state:6}
$n = N_D = \qty{1.0e16}{cm^{-3}}$; $p = n_i^2/n = \num{1.0e20}/\num{1.0e16} = \qty{1.0e4}{cm^{-3}}$.
\end{solution}

\begin{solution}{exo:b3:solid-state:7}
$E_F - E_i = kT\ln(n/n_i) = \qty{0.02585}{eV} \times \ln(10^6) = \qty{0.36}{eV}$.
\end{solution}

\begin{solution}{exo:b3:solid-state:8}
$kT = \qty{0.06894}{eV}$ at \qty{800}{K}; $n/N = \eu^{-2.3/0.1379} = \num{5.7e-8}$; in a mole $\num{6.02e23} \times \num{5.7e-8} = \num{3.4e16}$
pairs.
\end{solution}

\begin{solution}{exo:b3:solid-state:9}
The mass of a cell is $\rho a^3 = \qty{5.70}{g.cm^{-3}} \times (\qty{4.300e-8}{cm})^3 = \qty{4.532e-22}{g}$, that is
$\qty{4.532e-22}{g} \times N_A = \qty{272.9}{g}$ per mole of cells, or \qty{68.23}{g} per oxygen. Then
$(1 - x) \times 55.8 + 16.0 = 68.23$ gives $1 - x = 0.936$: $x = 0.064$, $\mathrm{Fe_{0.936}O}$.
\end{solution}

\begin{solution}{exo:b3:solid-state:10}
$\ln(\sigma T)$: $-3.51$, $-1.30$, $0.36$ at $1/T = 3.333$, $2.857$, $\qty{2.500e-3}{K^{-1}}$; slope between the end points
$(0.36 + 3.51)/(-\qty{8.33e-4}{K^{-1}}) = -\qty{4.65e3}{K}$ (the middle point lies on the line), so
$E_a = \qty{4.65e3}{K} \times \qty{8.617e-5}{eV.K^{-1}} = \qty{0.40}{eV}$.
\end{solution}

\begin{solution}{exo:b3:solid-state:11}
$\tfrac12\ce{O2} \rightleftharpoons \mathrm{O_O^{\times} + V_M' + h^{\bullet}}$, $K = [\mathrm{V_M'}][\mathrm h^\bullet]/p(\ce{O2})^{1/2}$. Charge balance $[\mathrm h^\bullet] = [\mathrm{V_M'}]$, so
$[\mathrm h^\bullet]^2 = K\,p(\ce{O2})^{1/2}$ and $[\mathrm h^\bullet] \propto p(\ce{O2})^{1/4}$; the conductivity follows the \omterm{def:b3:solid-state:carriers}{holes}.
\end{solution}

\begin{solution}{exo:b3:solid-state:12}
$n/N = \sqrt{N_i/N}\,\eu^{-\Delta H_F/2kT} = \sqrt 2\,\eu^{-1.4/(2 \times 0.05170)} = 1.414 \times \num{1.32e-6} = \num{1.9e-6}$.
\end{solution}

\begin{solution}{pb:b3:solid-state:1}
\textbf{1.} $\ce{Y2O3} \xrightarrow{\ce{ZrO2}} \mathrm{2\,Y_{Zr}' + 3\,O_O^{\times} + V_O^{\bullet\bullet}}$.

\textbf{2.} Sites: two cation sites and four oxygen sites (three filled, one empty),
the 1\,:\,2 ratio of $\ce{ZrO2}$. Mass: 2 Y and 3 O on each side. Charge: $2 \times (-1) + 2 = 0$.

\textbf{3.} One vacancy for two yttrium ions: one half per yttrium.

\textbf{4.} Per $(\ce{ZrO2})_{0.92}(\ce{Y2O3})_{0.08}$ there are 0.92 Zr and 0.16 Y, 1.08 cations: $x = 0.16/1.08 = 0.148$,
$\mathrm{Zr_{0.852}Y_{0.148}O_{1.926}}$.

\textbf{5.} $(x/2)/2 = 0.0370$: 3.7\,\% of the oxygen sites are empty.

\textbf{6.} $8 \times 0.0370 = 0.296$ vacancies per cell.

\textbf{7.} $a^3 = (\qty{5.14e-8}{cm})^3 = \qty{1.358e-22}{cm^3}$: $0.296/\qty{1.358e-22}{cm^3} = \qty{2.2e21}{cm^{-3}}$.

\textbf{8.} Each oxide ion must hop over a barrier of \qty{1.0}{eV} into a neighbouring
vacancy; the fraction of attempts that succeed, $\eu^{-E_a/kT}$, grows very fast with $T$.

\textbf{9.} At \qty{1073.15}{K}, $kT = \qty{0.09248}{eV}$: $A = \sigma T\eu^{E_a/kT} = 0.030 \times 1073.15 \times \eu^{10.81} = \qty{1.6e6}{S.K.cm^{-1}}$.

\textbf{10.} At \qty{973.15}{K}: $\sigma = (A/T)\eu^{-11.92} = \qty{0.011}{S.cm^{-1}}$.

\textbf{11.} At \qty{573.15}{K}: $\sigma = (A/T)\eu^{-20.25} = \qty{4.5e-6}{S.cm^{-1}}$.

\textbf{12.} $R = L/(\sigma S) = \qty{0.10}{cm}/(\sigma \times \qty{0.20}{cm^2})$: \qty{46}{\ohm} at \qty{700}{\celsius}, $\qty{1.1e5}{\ohm}$ at \qty{300}{\celsius}.

\textbf{13.} $R < \qty{1}{k\ohm}$ needs $\sigma > \qty{5.0e-4}{S.cm^{-1}}$; solving $(A/T)\eu^{-E_a/kT} = \num{5.0e-4}$ (by trial)
gives $T \approx \qty{760}{K}$, about \qty{490}{\celsius}.

\textbf{14.} Exhaust is cold after a start and at low load; the heater brings the
electrolyte to its working temperature within seconds, so that the control
works from the start, when most pollutants are emitted.

\textbf{15.} $\ce{O2 + 4e- -> 2O^2-}$; in \omterm{def:b3:solid-state:kroger-vink}{Kröger--Vink notation} $\ce{O2} + \mathrm{2\,V_O^{\bullet\bullet} + 4\,e'} \longrightarrow \mathrm{2\,O_O^{\times}}$.

\textbf{16.} Oxygen is reduced on the air side and the oxide ions are oxidised back
to oxygen on the exhaust side: the cell transfers $\ce{O2}$ from air to exhaust, with
$\Delta_rG = RT\ln(p_{\mathrm{exh}}/p_{\mathrm{air}})$ per mole of $\ce{O2}$ and four electrons: $E = -\Delta_rG/4F = (RT/4F)\ln(p_{\mathrm{air}}/p_{\mathrm{exh}})$.

\textbf{17.} $RT/4F = 8.3145 \times 973.15/(4 \times 96485) = \qty{0.02096}{V}$.

\textbf{18.} Zero: the two pressures are equal.

\textbf{19.} $(RT/4F)\ln 10 = \qty{48.3}{mV}$ per decade.

\textbf{20.} $\qty{0.02096}{V} \times \ln(0.2095/0.010) = \qty{0.02096}{V} \times 3.04 = \qty{0.064}{V}$.

\textbf{21.} $\qty{0.02096}{V} \times \ln(0.2095/\num{1.0e-20}) = \qty{0.02096}{V} \times 44.49 = \qty{0.93}{V}$.

\textbf{22.} Rich exhaust contains carbon monoxide, hydrogen and unburnt
hydrocarbons; on the platinum electrode the little oxygen left reacts with them,
and $p(\ce{O2})$ is set by the equilibria $\ce{2CO + O2 <=> 2CO2}$ and $\ce{2H2 + O2 <=> 2H2O}$, which leave a tiny
value at \qty{700}{\celsius}.

\textbf{23.} $p = 0.2095\,\eu^{-0.45/0.02096} = \qty{1.0e-10}{bar}$.

\textbf{24.} Around the stoichiometric ratio, the exhaust passes from a slight excess
of oxygen to a slight excess of CO and hydrogen: $p(\ce{O2})$ falls by some eighteen
powers of ten for a small change in fuel, and the voltage jumps from below
\qty{0.1}{V} to about \qty{0.9}{V}. The engine control adds fuel when the voltage is low and
removes it when high, holding the mixture at the stoichiometric point where the
three-way catalyst converts CO, hydrocarbons and nitrogen oxides at once.

\textbf{25.} In rich exhaust at \qty{700}{\celsius} the sensor gives about \qty{0.93}{V}.
\end{solution}
